\documentclass[12pt]{article}


\usepackage[english,russian]{babel}

\usepackage{spbu-ff-3kurs}

\usepackage{amsmath}
\usepackage{graphicx}
\usepackage[colorlinks=true, allcolors=blue]{hyperref}

\begin{document}

\FirstPage{СВ.5011.2021 «Физика»}{21.Б08-фз}{Иванова Инна Сергеевна}{Фамилия Имя Отчество}{должность кафедры ФВЭиЭЧ }{2024}


\ToC


\section*{Введение}
\addcontentsline{toc}{section}{Введение}
В данной работе рассматривается решение уравнения Фока-Клейна-Гордона (уравнение ФКГ) в пространстве постоянной кривизны де Ситтера. Работаем со свободным массивным полем. Решение найдено в виде свертки координаты пространства де Ситтера X со свето-подобным вектором нулевой длины V в некоторой степени. 

Рассматриваемые координатный вектор и вектор V будут переписаны с помощью спинорного формализма для удобства учета лоренц-ковариантной формы. Будет описан математический аппарат работы с подобного рода объекты.

Планируется на основе результатов этой задачи построить голографическое описание пространства де Ситтера, аналогичное AdS/CFT соответствию. Для этого надо разработать соответствующую теорию поля, а именно найти явный вид операторов рождения-уничтожения. 


\section*{Основные результаты}
\addcontentsline{toc}{section}{Основные результаты}
От описания поля некоторой функцией переходим к разложению по полному ортогональному набору. Был произведен аналог Фурье разложения, но в криволинейном пространстве. Как известно, коэффициентами разложения в квантовой теории поля являются операторы рождения и уничтожения.

 Операторы поля рассматриваются на V-пространстве векторов нулевой длины, где V можно представить через Вейлевский спинор. Полученные операторы рождения-уничтожения интерпретируются как полевые операторы на спинорном пространстве и использованы для поиска дуальной теории.
 

\section*{Заключение}
\addcontentsline{toc}{section}{Заключение}
В результате проделанной работы удалось решить уравнение ФКГ в нужном виде и выразить решение через параметры задачи, а также переписать эту конструкцию в терминах спиноров. Не до конца удалось разобраться с дуальной теорией, аналогичной AdS/CFT соответствию, работа над этим будет продолжена позже\cite{V04r,Asch_self,Hwa,Forster_Nelson77}.


\newpage
\addcontentsline{toc}{section}{Список использованных литературных источников  и информационных материалов}
\bibliographystyle{unsrt}
\bibliography{sample}


\section*{Перечень использованного оборудования, в том числе оборудования  Научного парка СПбГУ}
\addcontentsline{toc}{section}{Перечень использованного оборудования, в том числе оборудования  Научного парка СПбГУ }


\section*{Приложения}
\addcontentsline{toc}{section}{Приложения}


\end{document}